\section{引言}

区域风电集群预测直接服务于电力系统调度、备用容量配置和新能源消纳。与单场站预测相比，区域集群预测关注多个风电场站的聚合出力，既受区域天气过程驱动，也受预测起点前场群真实运行状态、场站间同步变化和短期功率惯性影响。随着网格化数值天气预报和深度气象预报模型的发展，未来气象预报已经成为日内多步风电预测中的关键外生输入。

对于调度场景而言，区域尺度预测更接近实际功率平衡需求。单个场站的局部误差在聚合后可能相互抵消，但区域尺度的天气传播、同步爬坡和限电恢复等现象仍会造成显著预测偏差。因此，区域集群预测不仅需要描述未来风场变化，还需要在预测起点处对当前场群运行状态进行校准。

本文关注业务化区域风电集群 24 h 预测。该任务具有明确的信息可用性约束：未来网格化气象预报覆盖整个目标窗口，是预测未来功率变化的主要驱动；而场站级实测功率只在目标窗口前可用，主要表征当前出力水平、短期惯性和多场站耦合状态。仅使用未来气象的 Weather-only 方法能够捕捉主要天气驱动，但会忽略预测起点前的真实运行状态，尤其容易在预测窗口前段产生状态失配误差。

这一状态失配并不能简单通过增加气象编码器复杂度解决。即使气象预报能够较好描述未来风场，气象到功率的映射仍受当前功率水平、场站运行状态和短时惯性的影响。如果模型只从未来气象场直接回归区域功率，则预测窗口开始阶段容易偏离真实起点；如果将历史功率与未来气象对称拼接，又可能弱化二者在业务时间轴上的不同角色。

已有风电预测研究从风速时序深度模型 \cite{ref1}、网格化空间表示 \cite{ref2}、场群时空关联建模 \cite{ref3,ref4,ref5,ref6,ref7,ref8,ref9,ref10,ref11,ref12}、NWP 校正和多模态融合 \cite{ref13,ref14,ref15,ref16,ref17,ref18} 等角度提升预测性能。然而，许多融合框架仍将历史功率和未来气象作为相对对称的输入，通过拼接、并行编码或特征级注意力进行交互。本文认为，在业务化 24 h 区域风电预测中，更自然的建模方式是非对称气象-功率融合：未来气象作为主驱动记忆，近期功率作为预测起点状态条件，对未来气象信息的读取过程进行约束。

基于这一思想，本文提出一种运行状态约束的气象-功率融合预测网络。模型首先将后报气象和未来气象预报编码为可查询的气象记忆；随后通过多尺度图编码器从最近场站级功率中提取当前场群状态，并生成动态功率上下文 token；最后将这些 token 前置注入气象记忆，使每个 future query 在 cross-attention 中同时读取气象主驱动和功率状态条件，输出 24 步区域集群预测序列。

本文主要贡献如下：
\begin{itemize}
\item 从业务时间可用性出发，将区域风电集群预测表述为非对称气象-功率融合问题，明确未来气象的主驱动作用和近期功率的运行状态约束作用。
\item 提出功率状态条件化的气象-功率融合网络，在 memory 层面注入动态功率上下文 token，而不是仅在输出层进行特征拼接。
\item 在吉林区域 58 个风电场站数据上验证方法有效性。结果表明，本文方法整体优于 Weather-only 和多种网格气象编码基线；逐预测步分析和消融实验说明，真实功率状态能够增强 24 h 预测窗口内的状态适配能力，且起始时段收益尤为显著。
\end{itemize}

\figwide{0.78\textwidth}{../../paper_image_use/method/Fig01_svg.pdf}{业务时间协议。}{近期场站功率仅在目标窗口前可用，未来气象预报覆盖完整预测窗口；本文在该时间可用性约束下进行气象-功率融合预测。}

\figwidelabel{0.78\textwidth}{../../paper_image_use/method/Fig02_svg.pdf}{整体框架。}{未来气象构成主记忆，近期场站功率生成动态上下文 token，并在 memory 层面对气象读取过程进行状态条件化。}{fig:framework}
